\chapter{Growth Operators and the Algebra $\Dop$}\label{ch:growth}

\begin{workplan}{\Recast}
\wpgoal{Keep the typed structural operators. Rebuild the algebra $\Dop$ on
\emph{measured} relations rather than asserted ones, and be honest that in V2
this chapter is mostly a programme rather than a result.}

\wpsource{This chapter's existing text;
\texttt{MATH\_TOOLBOX.md}~I.1 (typed ops replacing Pachner moves), I.7 (the
curvature controller --- derived, unit-tested, never wired);
\texttt{DOCS/Phase I --- Understand the Geometry o.txt} (the phased programme:
invariants, then primitives, then interaction laws, then the algebra).}

\wpsteps{
  \item \textbf{Keep the typed operators} $\opbind$, $\opcollapse$,
    $\opsplit$, $\opmerge$, $\opcone$, each logged to a mutation history.
    Keep the reason they replaced Pachner moves: a knowledge graph is not a
    manifold, so the preservation guarantee does not transfer. \Keep.
  \item \textbf{Add the invariants step \emph{before} the primitives.} The V0
    text hunts for primitive operations first. The correct order --- and the
    one the source programme itself corrects to --- is invariants first:
    ``primitive'' is only meaningful relative to what must not move. Each
    operator must now declare its effect on $\mathrm{K}(\mathbb{K})$
    (Chapter~\ref{ch:identity}). Make that declaration part of the definition
    of an operator, not a later audit.
  \item \textbf{Derive the relations; do not invent them.} The commutators are
    to be \emph{measured}: apply candidate primitives in both orders on small
    explicit complexes and compute the difference. State clearly that
    $\Dop = T(V)/I$ is a target, and that $I$ is empty until such measurements
    exist. An algebra asserted without its relations is bookkeeping.
  \item \textbf{Report the independence-relation test honestly.} If the
    independence relation is empty, $\Dop$ is the free algebra and the whole
    apparatus counts directed acyclic graphs. Record this as the falsifier for
    this chapter, and record that it has not been run for V2's operator set.
  \item \textbf{Keep the curvature controller with its warning.} ``Run the
    flow, skip the surgery'' is right, and Forman's weights are \emph{not}
    free --- they are fixed by whichever Laplacian the system actually prints
    $\omega$ and $\rho$ from, so using a Hebbian coupling as the Forman weight
    computes the curvature of a \emph{different} operator: a consistent-looking
    number about the wrong thing. Keep the flagged prediction that a
    low-precision edge may drift positive and be folded \emph{because we are
    unsure of it}, which is backwards; check the sign behaviour before
    trusting the controller.
  \item \textbf{Keep the floor derivation as a template.} The exact
    multiplicative flow can never sever an edge; what severs is the
    forward-Euler discretisation flipping sign --- so the old weight floor was
    patching a numerical artefact that had been misdiagnosed as a safety
    mechanism. The replacement is exponential integration plus a control
    barrier, which turns out to be exactly the object neuroscience calls
    weight-dependent plasticity. Keep also the floating-point trap: bare
    exponential integration underflows to zero and severs silently; only the
    barrier survives floating point.
  \item \textbf{State the implementation status plainly.} No line of the
    curvature controller has ever moved a real edge weight. \Proposed.
}

\wpdone{Every operator declares its effect on the invariant, and the chapter
distinguishes what is proved from what is a programme awaiting measurement.}
\end{workplan}

\begin{intu}\Intu
Learning \emph{is} change of shape. The complex binds organs that work
together, prunes idle ones, and --- when a procedure repeatedly succeeds
--- promotes it to a first-class operator. Four axes of growth, one of
which grows the very set of moves the system can make.
\end{intu}

\section{The structural operators}

\begin{definition}[Base moves]\label{def:ops}
The base moves on $\CC$ are the typed operators
\[
  \{\opbind,\ \opcollapse,\ \opsplit,\ \opmerge\},
\]
each logged to an immutable mutation history.
\end{definition}

\begin{remark}[``Pachner move'' was discarded, correctly]
Earlier drafts called these Pachner moves. A knowledge complex is not a
manifold, so the invariance guarantees that make Pachner moves useful
--- that any two triangulations of a PL manifold are connected by a
finite sequence of them --- simply do not transfer. Borrowing the name
would have imported a theorem that does not apply.
\end{remark}

\begin{construction}[The four growth axes]\label{con:axes}
Evolution proceeds along
\[
  \underbrace{\delta\CC_{\text{fine}}}_{\text{concepts}},\quad
  \underbrace{\delta\CC_{\text{coarse}}}_{\text{coalitions}},\quad
  \underbrace{\delta\Fs}_{\text{stalk and restriction data}},\quad
  \underbrace{\delta\Dop}_{\text{new operators}} .
\]
The fourth axis is the essential one: the operator algebra is
time-dependent, $\Dop(t)\to\Dop(t+1)$.
\end{construction}

\begin{definition}[Hebbian coupling and decay]\label{def:hebb}
For a coalition $\sigma$ with weight $w(\sigma,t)\in[0,1]$:
co-activation increments $w$; idle steps apply
$w(\sigma,t+1)=(1-\gamma)w(\sigma,t)$ for a decay rate $\gamma\in(0,1)$;
a coalition falling below a bind threshold \emph{collapses}. A pair past
the bind threshold is a $1$-simplex of $\Kc$.
\end{definition}

\begin{proposition}[Idle coalitions collapse in finite time]\label{prop:decay}
\Proved\ If $\sigma$ is idle from time $t_0$ and the bind threshold is
$\theta>0$, then $\sigma$ collapses at the first
$t\ge t_0+\log_{1/(1-\gamma)}\bigl(w(\sigma,t_0)/\theta\bigr)$.
\end{proposition}
\begin{proof}
Iterating the decay gives $w(\sigma,t_0+m)=(1-\gamma)^mw(\sigma,t_0)$,
which falls below $\theta$ as soon as
$m\ln(1-\gamma)<\ln\bigl(\theta/w(\sigma,t_0)\bigr)$; since
$\ln(1-\gamma)<0$ this is $m>\log_{1/(1-\gamma)}(w(\sigma,t_0)/\theta)$.
\end{proof}

\begin{construction}[Consolidation and microneurons]\label{con:micro}
Over the weight filtration, persistent homology reads which coalitions
and concept-cycles \emph{survive}; survivors are promoted into permanent
structure. A \emph{microneuron} is a promoted sub-tree: when a composite
DAG repeatedly resolves conflict \emph{and} passes verification, it is
named and thereafter emitted by the planner as a single atomic generator
of $\Dop$. This is $\delta\Dop$ made concrete. The promotion gate is
\textsc{VerifyOp} (Chapter~\ref{ch:verify}): a composite becomes a
generator only if it survives the external oracle.
\end{construction}

\begin{hon}\Hon
The requirement that promotion pass \textsc{VerifyOp} is doing real
work. Without it, ``repeatedly resolved conflict'' promotes whatever
procedure most reliably makes the organs agree --- which, in a system
whose only cheap signal is agreement, selects for procedures that
manufacture consensus rather than truth. Coherence is a control signal,
not a fitness function, and the gate is what keeps that distinction
enforced.
\end{hon}

\chapter{The Fine Dynamics}\label{ch:flow}

\begin{intu}\Intu
Inside an organ, reasoning is a \emph{flow}. New material perturbs the
internal state and a gradient descent pulls it back toward internal
coherence, subject to hard constraints (axioms that must hold) and soft
constraints (inferred facts, bendable).
\end{intu}

\begin{construction}[Fine state and deductive flow]\label{con:flow}
The fine state of an organ is a configuration $x\in(\Rd)^m$ over its $m$
active concepts, with a scalar \emph{conflict functional} $\Phi(x)\ge0$
measuring internal inconsistency --- the fine-level analogue of
$\discord$. Reasoning integrates
\begin{equation}\label{eq:flow}
  \dot x=-\nabla\Phi(x)
\end{equation}
by a fixed-step fourth-order Runge--Kutta scheme. The reported
\emph{$\delta$-strain} of a step is the achieved change $\Delta\Phi$.
Two constraint types shape $\Phi$: \emph{rigid} constraints
(user-stated axioms) act as hard orthogonal projections onto an affine
subspace via a projector $P$; \emph{fluid} constraints (inferred facts)
act as soft quadratic penalties.
\end{construction}

\begin{proposition}[The flow is a descent]\label{prop:descent}
\Proved\ Along any solution of \eqref{eq:flow},
$\frac{d}{dt}\Phi(x(t))=-\norm{\nabla\Phi(x(t))}^2\le0$, with equality
exactly at critical points. If $\Phi$ is bounded below --- which it is,
by $\Phi\ge0$ --- then $\Phi(x(t))$ converges and
$\nabla\Phi(x(t))\to0$ along a subsequence.
\end{proposition}
\begin{proof}
Chain rule: $\frac{d}{dt}\Phi=\inner{\nabla\Phi}{\dot x}
=-\inner{\nabla\Phi}{\nabla\Phi}$. Monotone and bounded below implies
convergent; integrating
$\int_0^\infty\norm{\nabla\Phi}^2\,dt=\Phi(x(0))-\lim\Phi<\infty$
forces $\liminf\norm{\nabla\Phi}=0$.
\end{proof}

\begin{remark}[Observed sign behaviour]
On a live run, injecting a new concept raised $\Phi$ by a structural
jump ($\delta$-strain $=+0.202$) and the subsequent flow pulled it back
down ($\delta$-strain $=-0.081$): a structural expansion followed by a
deductive relaxation, which is the intended dynamics. Note that the
expansion is \emph{not} a descent step --- it is an exogenous change of
the functional itself, and only the relaxation is governed by
Proposition~\ref{prop:descent}.
\end{remark}

\begin{hon}\Hon
The closed form of $\Phi$ is an implementation-level choice --- a
penalised quadratic on the fine complex --- and is deliberately not
over-specified here. What is mathematically committed is \eqref{eq:flow}
(a gradient flow on a nonnegative conflict functional) and the
rigid/fluid split as hard projection versus soft penalty. Convergence to
a \emph{global} minimum is not claimed and would be false for a
non-convex $\Phi$.
\end{hon}

\chapter{Scheduling: Antichains, and Why ``Operad'' Is Not Earned}\label{ch:operad}

\begin{intu}\Intu
A reasoning act is a small program: a directed acyclic graph of
operators, some depending on others. Independent operators can run at
once, and the scheduler slices the DAG into layers of mutually
independent work.
\end{intu}

\begin{definition}[Antichain foliation]\label{def:foliation}
A reasoning act is a DAG $T=(N,\prec)$ of operator nodes. The scheduler
foliates $T$ into the antichain layers of $\prec$: layer $0$ is the set
of minimal nodes, and layer $\ell+1$ is the set of nodes all of whose
predecessors lie in layers $\le\ell$. Each layer is executed in
parallel. Correctness follows because nodes in one antichain share no
$\prec$-dependency. Node \emph{support} sets index which fine-complex
vertices an operator reads and writes, so slices with disjoint support
are safe to co-execute.
\end{definition}

\begin{remark}[The intended structure is a trace monoid]\label{rem:tracemonoid}
What Definition~\ref{def:foliation} describes is antichain layering of a
DAG. An operad requires multi-input composition together with
equivariance, associativity and unit axioms, none of which are stated or
used anywhere. The structure the design is reaching for is a
\emph{trace monoid} $\mathbb M(\Sigma,I)$ in the sense of Mazurkiewicz:
the free monoid on the operator alphabet modulo $ab=ba$ for pairs in an
\emph{independence relation} $I$, whose canonical Foata normal form is
exactly the antichain layering the scheduler computes. That
identification would be a \emph{sharpening} rather than a downgrade,
since it supplies normal forms, a Hilbert series, and a canonical form
for deduplication. The terminology should be corrected regardless of
whether the identification is completed.
\end{remark}

\begin{hon}\Hon
\Refuted\ The trace-monoid identification \textbf{does not currently
hold}, and the reason is three separate defects, measured 2026-07-29
against the deployed operator set.

\emph{(i) The relation is not symmetric.} In the scheduler, admitting
\emph{any} empty-support node --- including a read-only one --- marks
the whole slice as a global mutation and blocks every later node. So $8$
of the $15$ operator pairs can share a slice in one order and not in the
other. An independence relation is symmetric by definition; a
non-symmetric relation defines no trace monoid, hence no Foata normal
form and no Hilbert series. This is very likely an unintended scheduling
behaviour and is worth fixing on its own merits.

\emph{(ii) The support sets are placeholders.} Every support accessor
returns a hardcoded literal --- $\varnothing$, $\{0\}$ or $\{1\}$ ---
with three of the six operators returning the same $\{0\}$, and a
comment standing in for an analysis of what the operator actually
touches. Any invariant computed from these describes six constants, not
the system.

\emph{(iii) The convention for $\varnothing$ is inverted.} In the engine
an empty support means \textsc{global}: touches everything, must run
isolated. Read as a set, $\varnothing$ is disjoint from everything and
would mean independent of everything. The two readings give capacities
differing by a factor of $1.8$ ($\kappa=3.00$ versus $\kappa=5.45$,
against $\kappa=6$ for the free monoid), so the quantity is not well
defined until the convention is fixed.

Under the reading the engine actually implements, the algebra is
\emph{nearly free} --- about $9\%$ below the free monoid --- because
three of six operators share one support. The order of repair is: fix
the symmetry defect; derive real read/write sets, including the shared
state the current model ignores entirely (the store, the obstruction
counter $\Omega$, the mutation log, and the coarse stalks that $\pi_v$
overwrites after every execution); only then recompute $I$. Operators
that all mutate $\Omega$ do not commute whatever their vertex supports
say. Until that is done, no capacity, Koszulity, or resolution claim
about $\Dop$ should appear anywhere.
\end{hon}

\chapter{The Lift--Project Adjunction}\label{ch:adjunction}

\begin{workplan}{\Keep}
\wpgoal{Keep the adjunction --- it is standard, genuinely proved, and now
does double duty as the memory architecture and as the candidate reproduction
morphism.}

\wpsource{This chapter's existing text;
\texttt{MATH\_TOOLBOX.md}~II.2 (the $\mathbb{K}/W$ two-complex model, the
consolidation rule, the measured $Q(t)$ trajectory), II.1 (what was tried and
abandoned for memory identity first).}

\wpsteps{
  \item \textbf{Keep the theorem.} For a full subposet inclusion $\iota$:
    $\iota^*$ restricts (instantiate, store $\to$ cache), $\iota_!$ is
    extension by zero (write back only what was touched), $\iota_*$ fills
    untouched cells by limits, and the (co)units are isomorphisms so
    $\iota^*\iota_! \cong \Id$. \Proved, standard.
  \item \textbf{State the write-back choice and its reason.} $\iota_!$, never
    $\iota_*$: the cache must assert \emph{nothing} about untouched cells,
    which is the sheaf-theoretic form of ``no measurement is not a measurement
    of zero.''
  \item \textbf{Present consolidation as the anti-interference mechanism.}
    $\gamma_0 \ll 1$ is what lets new learning coexist with old memory without
    overwriting it, and the session's state $s_W$ is \emph{never} written
    back --- content and wiring persist, the episode does not.
  \item \textbf{Report that $\gamma_0 \ll 1$ was forced twice, independently.}
    Once by the anti-interference argument, and once by measurement: no
    restriction-map parameterisation is affordable to learn from a single tick
    (the cheapest is over-parameterised threefold, a full orthogonal map by
    over five hundred). Two independent routes to one conclusion is worth
    stating as such.
  \item \textbf{Keep $Q(t)$ as an instrument.} Mean $\norm{R - \Id}_F^2$;
    constant sheaf $\iff Q = 0$; measured moving monotonically off zero across
    four ticks, and \emph{exactly} zero when nothing is verified. ``$Q(t)$
    leaving zero is not a proxy for consolidation; it is the definition.''
    \Measured.
  \item \textbf{Keep the two bugs the tests caught}, since both would have
    rotted silently: a constructor that pre-filled every edge with identity so
    learned maps were dropped on rebuild, and a single Householder reflection
    landing in the wrong component of $O(d)$ so consolidation correctly refused
    to interpolate and $Q(t)$ would have stayed zero with every test passing.
  \item \textbf{Add the forward reference to reproduction.} This same adjoint
    pair is Chapter~\ref{ch:reproduction}'s candidate morphism. Note it here.
  \item \textbf{Keep the abandoned alternatives with their reasons.} Formal
    Concept Analysis \Refuted\ as substrate (data-determined; its canonicity
    comes precisely from discarding the path-dependence that makes a memory a
    memory --- memory must \emph{grow}, not be recomputed), and the
    $K_0$/Jordan--Hölder schema \Refuted\ (both candidate categories fail for
    opposite reasons). Both were dissolved by the growth-history identity this
    chapter provides.
}

\wpdone{The adjunction is proved, the write-back choice is justified by the
$\bot$-invariant, and its second role in reproduction is signposted.}
\end{workplan}

\begin{intu}\Intu
Two worlds: free-form natural language, and structured reasoning DAGs.
One map \emph{lifts} language into a validated DAG; a partner map
\emph{projects} a DAG back into language. Meaning that carries no
reasoning content collapses under the lift --- it lands in the kernel.
\end{intu}

\begin{construction}[The boundary morphisms]\label{con:adjunction}
Let $\mathcal L$ send a natural-language request to a validated
reasoning DAG, and $\mathcal M$ send a DAG to a natural-language
rendering. The lift is guarded by a validator rejecting malformed DAGs
at the boundary: duplicate identifiers, dangling children, cycles by
depth-first colouring, absence of a terminal output node, and
disconnection. The downstream engine therefore only ever receives
well-formed programs.
\end{construction}

\begin{remark}[The kernel behaves as intended]
Semantically void input maps into $\ker\mathcal L$: a joke embedded in a
mathematical prompt produced zero DAG nodes, exactly as a kernel element
should. This is an observation on one input, not a theorem.
\end{remark}

\begin{hon}\Hon
$\mathcal L\dashv\mathcal M$ is the intended categorical \emph{framing}
--- lift left adjoint to project --- and it organises the design
usefully. The unit and counit are not constructed here and the triangle
identities are not verified, so ``adjunction'' should be read as a
structuring principle plus two concrete morphisms with an observed
kernel behaviour, and not as a proved theorem. Stating this is part of
the honesty invariant rather than an apology for it.
\end{hon}
